% chapters/10/rangkuman.tex
\bagian{Rangkuman}

\begin{rangkumanbox}
\begin{itemize}
  \item Peta berpetak diubah menjadi graf dengan tiga aturan: setiap petak
        kosong menjadi simpul, setiap pasangan petak bersebelahan menjadi sisi,
        dan biaya setiap sisi bernilai satu.
  \item Peta disimpan sebagai rangkaian baris huruf agar mudah disunting dan
        mudah diperiksa dengan mata. Program memeriksa panjang baris serta
        keberadaan titik awal dan titik tujuan sebelum peta dipakai.
  \item Pencarian melebar (BFS) memeriksa simpul menurut banyaknya langkah,
        sehingga tidak memadai apabila biaya sisinya berbeda-beda.
  \item Algoritma Dijkstra memeriksa simpul menurut biaya sejauh ini, sehingga
        jalur yang dihasilkan dijamin termurah.
  \item Algoritma A* menambahkan perkiraan sisa jarak ke tujuan pada biaya
        sejauh ini, yaitu $f = g + h$, sehingga pencariannya terarah. Dijkstra
        adalah A* dengan $h$ bernilai nol.
  \item Tepi pencarian berupa antrean prioritas dari modul \code{heapq}.
        Sebuah petak dapat masuk antrean lebih dari sekali, dan isi yang sudah
        usang dibuang dengan memeriksa biaya terkini.
  \item Heuristik harus tidak pernah melebihkan jarak agar jalur yang
        dihasilkan tetap terpendek. Jarak Manhattan dan Euclidean memenuhi
        syarat itu untuk gerak empat arah, tetapi jarak Manhattan tidak
        memenuhinya untuk gerak delapan arah.
  \item Pengukuran pada peta bab ini: Manhattan memeriksa 274 petak, Euclidean
        353 petak, dan Dijkstra 437 petak, sedangkan biaya jalurnya sama pada
        ketiganya, yaitu 50. Heuristik mempengaruhi kinerja pencarian, bukan
        mutu hasilnya.
  \item Jalur direkonstruksi dengan menelusuri tabel asal dari titik tujuan
        mundur ke titik awal, lalu membalik urutannya. Keadaan tujuan yang tidak
        dapat dicapai harus ditangani.
  \item Jalur yang dihasilkan diperiksa keabsahannya: tidak kosong, dimulai
        dari titik awal, berakhir di titik tujuan, bersambung, dan tidak
        melewati dinding.
  \item Visualisasi pencarian memperlihatkan petak yang telah diperiksa, petak
        yang menunggu diperiksa, petak yang sedang diperiksa, dan jalur yang
        ditemukan.
  \item Pencarian jalur pada bab ini belum terhubung dengan NPC. Penghubungan
        keduanya, beserta pengaturan kapan pencarian diulang, merupakan pokok
        Bab~XI.
\end{itemize}
\end{rangkumanbox}
